\documentclass{article}
\usepackage[T1]{fontenc}
\usepackage{lmodern}
\usepackage{graphicx}
\usepackage{amsmath,amssymb}
\usepackage[a4paper,margin=2cm]{geometry}
\usepackage[absolute,overlay]{textpos}
\setlength{\TPHorizModule}{1mm}
\setlength{\TPVertModule}{1mm}
\usepackage[indonesian]{babel}
\usepackage{hyperref}
\setlength{\emergencystretch}{2em}
% unit: Halaman 95

\begin{document}

% === Halaman 95 (python/native) ===

95

D. Hill Cipher

• Dikembangkan oleh Lester Hill (1929), termasuk ke dalam polygram cipher yang 

berbasis aljabar linier

• Menggunakan m buah persamaan linier 

• Untuk m = 3 (enkripsi setiap 3 huruf), sistem persamaan linernya menjadi:

    C1 = (k11 p1 + k12p2 + k13 p3) mod 26

    C2 = (k21 p1 + k22p2 + k23 p3) mod 26

    C3 = (k31 p1 + k32p2 + k33 p3) mod 26





                                                C = KP









































=





















3

2

1

33

32

31

23

22

21

13

12

11

3

2

1

p

p

p

k

k

k

k

k

k

k

k

k

C

C

C

\end{document}
